\documentclass[10pt,a4paper]{amsart}
\usepackage[margin=19mm]{geometry}
\usepackage{amsmath,amssymb,mathtools}
\usepackage[scr=rsfs,cal=euler]{mathalfa}
\usepackage{xfrac}
\usepackage[unicode,hidelinks,bookmarks=false]{hyperref}
\newcommand{\ie}{i.e.\ }
\newcommand{\cf}{cf.\ }
\newcommand{\resp}{respectively\ }
\newcommand{\Subsection}{\subsection}
\providecommand{\nobreakdash}{\nobreak\hbox{-}}
\setlength{\emergencystretch}{2em}
\multlinegap0pt
\usepackage{bm}
\usepackage{bbm}
\usepackage{dsfont}
\usepackage{xspace}
\usepackage{cancel}
\usepackage{mathtools}
\usepackage{amsthm}
\makeatletter
\@ifundefined{thm}{\newtheorem{thm}{Thm}}{}
\@ifundefined{lem}{\newtheorem{lem}[thm]{Lemma}}{}
\makeatother
\def\geq{\geqslant}
\def\leq{\leqslant}
\title[Long-Time Existence of Quasilinear Wave Equations Exterior t]{Long-Time Existence of Quasilinear Wave Equations Exterior to Star-shaped Obstacle in $2\mathbf{D}$}
\author{Lai Ning-An, Ren Cui, Xu Wei}
\date{Source: arXiv:2507.06897v1 (CC BY 4.0)}
\begin{document}
\maketitle

\noindent\emph{Source and licence.} Long-Time Existence of Quasilinear Wave Equations Exterior to Star-shaped Obstacle in $2\mathbf{D}$. Version arXiv:2507.06897v1 (CC BY 4.0), distributed under the Creative Commons Attribution 4.0 International licence, as recorded in the arXiv licence metadata for this exact version and shown on the abstract page https://arxiv.org/abs/2507.06897v1. Full rights notice: licenses/arxiv-2507.06897-CC-BY-4.0.txt.\par\smallskip

\noindent\textit{Editorial scope.} Counted unit: the Lem whose source label is \texttt{lem5} in the article (source0.tex lines 595--609, source label \texttt{lem5}), delivered together with the article's own complete original proof of it (source0.tex lines 610--733, one \texttt{proof} environment of 124 lines). No mathematics is written, repaired or re-derived: statement and proof are the article's own text line for line. Required context retained uncounted: Y38 (lines 584--586); Y39 (lines 588--590). Every definition, display and label of those lines is retained unchanged. Omitted from this package and not redistributed: 1--583, 587--587, 591--594, 734--1935. Nothing inside a retained excerpt is omitted or altered: every retained body is delivered exactly as the article's lines are written, so no omission receipt of any kind accompanies this package. Citations: the retained excerpts carry no citation command at all, so no bibliography entry of the article is transcribed and no reference is redistributed. Reference adaptation: none was needed and none was made; every reference inside the delivered unit resolves inside it, so no unresolved reference or citation is produced. Printed slips kept literal: no printed word, symbol, hypothesis or number of the counted statement or of its proof was altered. Layout and class: the article's own class is \texttt{amsart} with packages including tikz, amsmath, amssymb, xcolor, geometry, hyperref, mathrsfs; the wrapper is the lane's pinned \texttt{amsart} editorial wrapper at 10pt on A4, which restates the theorem-like environments the retained text needs and adds the article's own macro definitions that the retained text needs (\texttt{\textbackslash geq}, \texttt{\textbackslash leq}), with extra wrapper packages (bm, bbm, dsfont, xspace, cancel, mathtools). The delivered numbering is the unit's own. Assets: no retained span contains a figure environment, a graphics-inclusion command, a PDF-page inclusion, a table or a TikZ picture; no image file is copied into this package, the package declares no asset dependency and no image file is redistributed with it. Licence: version arXiv:2507.06897v1 of this article is distributed under the Creative Commons Attribution 4.0 International licence, as recorded in the arXiv licence metadata for this exact version and shown on the abstract page https://arxiv.org/abs/2507.06897v1. A full rights notice is delivered at licenses/arxiv-2507.06897-CC-BY-4.0.txt. Characters: every retained excerpt is pure ASCII, so the delivered engine, XeLaTeX with its default Latin Modern text font, reports no missing character.
\begin{align}\label{Y38}
h^{\alpha\beta}=h^{\beta\alpha},
\end{align}

\begin{align}\label{Y39}
|h|=\sum_{\alpha,\beta=0}^2|h^{\alpha\beta}|\ll1.
\end{align}

\begin{lem}\label{lem5}
Let $\mathcal {K}$ be bounded and star-shaped with respect to the origin.
Assume that the perturbation terms $h^{\alpha\beta}$ satisfy conditions \eqref{Y38} and \eqref{Y39}.
Then for any smooth function $u$ with $u|_{\partial\mathcal {K}}=0$ and $supp\{u(t,\cdot)\}\subseteq\{x:|x|\leq t-2\}$, we have
\begin{align}\label{Y29}
&\frac12\frac{d}{dt}\int_{\mathbb R^2\backslash\mathcal
{K}}\bigg(|t-r||\partial u|^2+r\Big(u_t+u_r+\frac
u{2r}\Big)^2+\frac{|\Omega
u|^2}r+\frac {u^2}{4r}\nonumber\\&+h^{0\beta}(2Su+u)u_{\beta}
-th^{\alpha\beta}u_{\alpha}u_{\beta}\bigg)dx
\nonumber\\&\leq\int_{\mathbb R^2\backslash\mathcal {K}}\Box_h u
\big((S+1/2)u\big)dx+C\int_{\mathbb R^2\backslash\mathcal {K}}|Sh||\partial u|^2dx
\nonumber\\&\quad+C\int_{\mathbb R^2\backslash\mathcal {K}}|\partial h||\partial u|\big(|Su|+|u|\big)dx.
\end{align}
\end{lem}

\begin{proof}
By multiplying $\Box u$ by $u_t$, we have
\begin{align}\label{Y3}
u_t\Box u&=\frac12\partial_t(u_t^2)-\nabla\cdot(u_t\nabla
u)+\nabla u\cdot\nabla u_t
\nonumber\\&=\frac12\partial_t(u_t^2+|\nabla
u|^2)-\nabla\cdot(u_t\nabla u),
\end{align}
which implies
\begin{align}\label{Y4}
(tu_t)\Box u=\frac12\partial_t(tu_t^2+t|\nabla
u|^2)-\frac12(u_t^2+|\nabla u|^2)-\nabla\cdot(tu_t\nabla u).
\end{align}
By multiplying $\Box u$ by $x\cdot\nabla u$, we have
\begin{align}\label{Y5}
(x\cdot\nabla u)\Box u&=\partial_t(u_tx\cdot\nabla
u)-\frac12x\cdot\nabla(u_t^2)-\nabla\cdot(\nabla ux\cdot\nabla u)
\nonumber\\&\quad+|\nabla u|^2+\frac12x\cdot\nabla(|\nabla u|^2)
\nonumber\\&=\partial_t(u_tx\cdot\nabla
u)+\frac12\nabla\cdot(x|\nabla
u|^2-xu_t^2)+u_t^2\nonumber\\&\quad-\nabla\cdot(\nabla ux\cdot\nabla
u).
\end{align}
By multiplying $\Box u$ by $u/2$, we have
\begin{align}\label{Y6}
(u/2)\Box
u&=\frac12\partial_t(uu_t)-\frac12u_t^2-\frac12\nabla\cdot(u\nabla
u)+\frac12|\nabla u|^2.
\end{align}
The combination of equalities \eqref{Y4}, \eqref{Y5} and \eqref{Y6}
shows that
\begin{align}\label{Y7}
&\Box
u(tu_t+ru_r+u/2)\nonumber\\&=\frac12\partial_t(tu_t^2+t|\nabla
u|^2+2ru_tu_r+uu_t)\nonumber\\&\quad+\frac12\nabla\cdot(x|\nabla
u|^2-xu_t^2)-\nabla\cdot(tu_t\nabla
u)\nonumber\\&\quad-\nabla\cdot(\nabla ux\cdot\nabla
u)-\frac12\nabla\cdot(u\nabla u).
\end{align}
\par Now let's calculate the perturbation terms. By using the symmetry conditions \eqref{Y38}, for any $\gamma=0,1,2$, we have
\begin{align*}
&\big(h^{\alpha\beta}\partial_{\alpha}\partial_{\beta}u\big)\partial_{\gamma}u
\\&=\partial_{\alpha}\big(h^{\alpha\beta}u_{\gamma}u_{\beta}\big)
-\big(\partial_{\alpha}h^{\alpha\beta}\big)u_{\gamma}u_{\beta}-h^{\alpha\beta}\partial_{\gamma}\partial_{\alpha}u\partial_{\beta}u
\\&=\partial_{\alpha}\big(h^{\alpha\beta}u_{\gamma}u_{\beta}\big)
-\big(\partial_{\alpha}h^{\alpha\beta}\big)u_{\gamma}u_{\beta}
-\frac12\partial_{\gamma}\big(h^{\alpha\beta}u_{\alpha}u_{\beta}\big)
+\frac12(\partial_{\gamma}h^{\alpha\beta})u_{\alpha}u_{\beta},
\end{align*}
which implies
\begin{align}\label{Y35}
&\big(h^{\alpha\beta}\partial_{\alpha}\partial_{\beta}u\big)\big(Su\big)
\nonumber\\&=\big(h^{\alpha\beta}\partial_{\alpha}\partial_{\beta}u\big)\big(x_{\gamma}\partial_{\gamma}u\big)
\nonumber\\&=\partial_{\alpha}\big(h^{\alpha\beta}(Su)u_{\beta}\big)
-\big(\partial_{\alpha}h^{\alpha\beta}\big)\big(Su\big)u_{\beta}
-\frac12\partial_{\gamma}\big(x_{\gamma}h^{\alpha\beta}u_{\alpha}u_{\beta}\big)
\nonumber\\&\quad+\frac12(Sh^{\alpha\beta})u_{\alpha}u_{\beta}+\frac12h^{\alpha\beta}u_{\alpha}u_{\beta}.
\end{align}
The combination of equality \eqref{Y35} and the following equality
\begin{align*}
&\big(h^{\alpha\beta}\partial_{\alpha}\partial_{\beta}u\big)\big(u/2\big)
\\&=\frac12\partial_{\alpha}\big(h^{\alpha\beta}uu_{\beta}\big)
-\frac12\big(\partial_{\alpha}h^{\alpha\beta}\big)uu_{\beta}
-\frac12h^{\alpha\beta}u_{\alpha}u_{\beta}
\end{align*}
implies that
\begin{align}\label{Y31}
&\big(h^{\alpha\beta}\partial_{\alpha}\partial_{\beta}u\big)\big((S+1/2)u\big)
\nonumber\\&=\partial_{\alpha}\big(h^{\alpha\beta}(Su+u/2)u_{\beta}\big)
-\big(\partial_{\alpha}h^{\alpha\beta}\big)\big(Su\big)u_{\beta}
-\frac12\partial_{\gamma}\big(x_{\gamma}h^{\alpha\beta}u_{\alpha}u_{\beta}\big)
\nonumber\\&\quad+\frac12(Sh^{\alpha\beta})u_{\alpha}u_{\beta}-\frac12\big(\partial_{\alpha}h^{\alpha\beta}\big)uu_{\beta}.
\end{align}
By equalities \eqref{Y7} and \eqref{Y31}, we have
\begin{align}\label{Y37}
&(\Box_h
u)(Su+u/2)\nonumber\\&=\frac12\partial_t(tu_t^2+t|\nabla
u|^2+2ru_tu_r+uu_t+h^{0\beta}(2Su+u)u_{\beta}-th^{\alpha\beta}u_{\alpha}u_{\beta})\nonumber\\&\quad+\frac12\nabla\cdot(x|\nabla
u|^2-xu_t^2)-\nabla\cdot(tu_t\nabla
u)-\nabla\cdot(\nabla ux\cdot\nabla
u)-\frac12\nabla\cdot(u\nabla u)\nonumber\\&\quad+
\partial_i\big(h^{i\beta}(Su+u/2)u_{\beta}\big)
-\frac12\nabla\cdot\big(xh^{\alpha\beta}u_{\alpha}u_{\beta}\big)
\nonumber\\&\quad-\big(\partial_{\alpha}h^{\alpha\beta}\big)\big(Su\big)u_{\beta}
+\frac12(Sh^{\alpha\beta})u_{\alpha}u_{\beta}-\frac12\big(\partial_{\alpha}h^{\alpha\beta}\big)uu_{\beta}.
\end{align}
Let $\vec{n}=(n_1,n_2)$ be the unit outward normal to $\mathcal
{K}$. The boundary condition $u|_{\partial\mathcal {K}}=0$ shows
that $u_t=0$ and $\frac{\partial u}{\partial x_i}=n_i\frac{\partial
u}{\partial n}$, $i=1,2$, on $\partial\mathcal {K}$. Since $\mathcal
{K}$ is star-shaped with respect to the origin, we know $\langle
x,\vec{n}\rangle|_{\partial\mathcal {K}}\geq0$, which implies
\begin{align*}
&-\frac12\int_{\partial\mathcal {K}}\langle x,\vec{n}\rangle|\nabla u|^2dS
+\int_{\partial\mathcal {K}}\frac{\partial u}{\partial n}\big(x\cdot\nabla u\big)dS
\\&\quad-\int_{\partial\mathcal {K}}n_ih^{il}\big(x\cdot\nabla u\big)\partial_ludS
+\frac12\int_{\partial\mathcal {K}}\langle x,\vec{n}\rangle h^{\alpha\beta}u_{\alpha}u_{\beta}dS
\\&=\int_{\partial\mathcal {K}}\Big(\frac{1-h^{il}n_in_l}2\Big)\langle x,\vec{n}\rangle\Big|\frac{\partial u}{\partial n}\Big|^2dS\\
&\geq0,
\end{align*}
where we used the size condition \eqref{Y39}.
Thus by integrating equality \eqref{Y37} with respect to $x$ on the
domain $\mathbb R^2\backslash\mathcal {K}$, we obtain
\begin{align}\label{Y28}
&\frac12\frac{d}{dt}\int_{\mathbb R^2\backslash\mathcal {K}}(tu_t^2+t|\nabla
u|^2+2ru_tu_r+uu_t\nonumber\\&\qquad\qquad\quad+h^{0\beta}(2Su+u)u_{\beta}-th^{\alpha\beta}u_{\alpha}u_{\beta})dx
\nonumber\\&\leq\int_{\mathbb R^2\backslash\mathcal {K}}(\Box_h
u)(Su+u/2)dx+C\int_{\mathbb R^2\backslash\mathcal {K}}|Sh||\partial u|^2dx
\nonumber\\&\quad+C\int_{\mathbb R^2\backslash\mathcal {K}}|\partial h||\partial u|\big(|Su|+|u|\big)dx.
\end{align}
A straight calculation shows that
\begin{align}\label{Y11}
&tu_t^2+t|\nabla u|^2+2ru_tu_r+uu_t\nonumber\\&=|t-r|(u_t^2+|\nabla
u|^2)+\frac{|\Omega u|^2}r\nonumber\\&\quad+r\Big(u_t+u_r+\frac
u{2r}\Big)^2-\frac {u^2}{4r}-uu_r,
\end{align}
and
\begin{align}\label{Y27}
\int_{\mathbb R^2\backslash\mathcal
{K}}uu_rdx=-\int_{\mathbb R^2\backslash\mathcal {K}}u^2/(2r)dx.
\end{align}
The combination of \eqref{Y28}, \eqref{Y11} and \eqref{Y27} implies
that inequality \eqref{Y29} holds.
\end{proof}

\end{document}
